% page 121 du fac-simile (image p-120 du scan)
% bloc 170.2 x 259.3 mm ; marge gauche 31.8 mm
\pgc{10.160mm}{0.000mm}{
\l{}
\l{}
\l{}
\l{}
\l{\cel{55}lll.}
\l{\sou{dinamiko}. La cienco pri la kalkulo di la movi di la korpi qui}
\l{\cel{3}submisesas a la efiko da la forci qui movas li. - DEFIRS.}
\l{}
\l{\sou{dinamismo}. (\sou{filoz.}) Sistemo en qua materio konsideresas kom facita}
\l{\cel{3}ek forci qui esas ipse efikiva. - DEFIRS.}
\l{}
\l{\sou{dinamisto}. Partisano di dinamismo.}
\l{}
\l{\sou{dinamito}. Explozivo ek mixuro de nitroglicerino e substanco inerta}
\l{\cel{3}(sablo quarcoza, greso polvo, e c.) - inventita da Alfred Nobel.}
\l{\cel{3}- DEFIRS.}
\l{}
\l{\sou{dinamometro}. Instrumento por mezurar la povo di motoro, la forteso}
\l{\cel{3}muskulala di la homi, di la animali. - DEFIRS.}
\l{}
\l{\sou{dinamometrio}. Evaluo e komparo di la fortesi, per dinamometro.}
\l{}
\l{\sou{dinastio}. Serio de suvereni, decendinta de un ancestro. - DEFIRS.}
\l{}
\l{\sou{dindo}. (\sou{zool.})\cel{2}Galinaceo di qua la kaudo, kurvatra abanike, esas}
\l{\cel{3}extensebla rotatre quale olta di pavono. - L.\cel{1}\sou{meleagris gallo-}}
\l{\cel{3}\sou{pavo}. - FI.}
\l{}
\l{\sou{dinear}. (\sou{netrans.}) Repastar, en la vespero (cirkum l9 o 2O kloki).}
\l{\cel{3}- DEF.}
\l{}
\l{\sou{dinosaurio}. (\sou{paleontol.})\cel{2}Klaso de repteri, ek nur formi desaparin-}
\l{\cel{3}ta, qui vivis sur la tersurfaco dum la epoki prehistoriala se-}
\l{\cel{3}kundara, ed inkluzanta : animali ek grandesi tre diversa, de}
\l{\cel{3}kelka centimetri til 25 metri de longeso. - DEF.}
\l{}
\l{\sou{dinoterio}. (\sou{zool.}) Genero de mamiferi rostroza, ye qua konsistas}
\l{\cel{3}la familio \textquotedbl{}dinoteridi\textquotedbl{}. - DEF.}
\l{}
\l{\sou{diocezo}. (\sou{eklezio katolika}). Distrikto di qua la judicieso reli-}
\l{\cel{3}giala resortisas ad episkopo od arkiepiskopo. - DEF.}
\l{}
\l{\sou{dioika}. (\sou{bot}.) Di qua la floruli e la florini esas sur stipi}
\l{\cel{3}diversa.}
\l{}
\l{\sou{dioptro}. La fenestreto tra qua onu regardas por vizar punto per}
\l{\cel{3}la instrumento di qua la nomo esas \textquotedbl{}alidado\textquotedbl{}. - DEIRS.}
\l{}
\l{\sou{dioptrio}. (\sou{optiko}). Mezur-unajo por la vitro-diski di la binokli.}
\l{\cel{3}- DEFIRS.}
\l{}
\l{\sou{dioptriko}. (\sou{optiko)}. Teorio di la vitro-diski optikala. - DEF.}
\l{}
\l{\sou{dioramo}. Pikturo granda-dimensiona, facita ek farbi diafana sur}
\l{\cel{3}telo sen bordi videbla, di qua onu variigas la lumizo por}
\l{\cel{3}prizentar diversa iluzioni optikala a la personi qui spektas}
\l{\cel{3}lu de loko obskura. - DEFIRS.}
\l{}
\l{\sou{diorito}. (\sou{geol.}) Roko fairatra, ek feldspato e verda amfibolo.-}
\l{\cel{3}DEFIRS.}
}
